\documentclass[11pt]{amsart}
\usepackage{amsmath,amssymb,amsthm,mathtools}
\usepackage[margin=1in]{geometry}
\usepackage{array}
\usepackage{hyperref}

\newtheorem{theorem}{Theorem}[section]
\newtheorem{proposition}[theorem]{Proposition}
\newtheorem{lemma}[theorem]{Lemma}
\newtheorem{corollary}[theorem]{Corollary}
\theoremstyle{definition}
\newtheorem{definition}[theorem]{Definition}
\theoremstyle{remark}
\newtheorem{remark}[theorem]{Remark}

\newcommand{\RH}{\mathrm{RH}}
\renewcommand{\Re}{\operatorname{Re}}
\renewcommand{\Im}{\operatorname{Im}}
\newcommand{\half}{\tfrac12}
\newcommand{\fhat}{\widehat f}
\newcommand{\ghat}{\widehat g}
\newcommand{\HK}{\mathcal H_K}
\newcommand{\HE}{H(E_\xi)}
\newcommand{\Fix}{\operatorname{Fix}}
\DeclareMathOperator{\Res}{Res}

\title[The Kre\u\i n Index of the Weil Form]{The localized Weil form and its Kre\u\i n--Pontryagin index:\\
localization, obstruction, and rigidity}
\author{David Alejandro Trejo Pizzo}
\thanks{Independent Researcher. \texttt{dtrejopizzo@gmail.com}}


\begin{document}
\nocite{bgConnes99,bgCC21,bgBognar,bgAzizov,bgBK}
\begin{abstract}
The Riemann Hypothesis is equivalent to the positivity of an indefinite (Kre\u\i n)
quadratic form: the localized Weil form $Q$, whose negative index $\kappa$ equals the
number of nontrivial zeros of $\zeta$ in one open half of the critical strip
(equivalently $2m$ for $m$ off-line Klein quartets). This paper assembles a
complete functional-analytic anatomy of $Q$ and of the realization problem ``carry
$\kappa$ by an operator.'' We work with an autonomous definition of $Q$ on even real test
functions through an entire transform, proving its convergence and symmetries directly,
and we show that the indefiniteness has a single explicit source: each off-line zero acts
through an \emph{off-axis evaluation functional} $f\mapsto\fhat(\gamma-ib)$, $b=\beta-\half
\ne0$, contributing a Lorentzian $u^2-v^2$, while every on-line zero contributes the
nonnegative square $\fhat(\gamma)^2$. From this localization we derive three structural
conclusions. \emph{(Obstruction.)} No positive-definite reproducing kernel, no fixed-sign
Cauchy kernel whose parameters lie on one side of a line, and no Hermitian form depending
smoothly and non-degenerately on the zero's real part can carry $\kappa$; we prove all
three as elementary but general theorems and illustrate the second with the modular-surface
resonance pairings, which are sign-definite for every configuration of zeros in the strip.
A genuine carrier must therefore have an indefinite kernel \emph{or} a degeneration at the
critical line. \emph{(Mechanism and count.)} The Weil form is bilinear, not sesquilinear;
the sesquilinear evaluation form is always positive, whereas the bilinear evaluation-square
form carries signature exactly $(1,1)$ at each off-axis point. The exact global count
$\kappa=2m$ (twice the number of off-line quartets) then rests on an independence lemma and the de Branges
Pontryagin-index theorem, valid in the finite-index regime. \emph{(Rigidity.)} Of the three
de Branges axioms, conjugation is automatic from the symmetries, point-evaluation continuity
is the reproducing-kernel hypothesis, and the difference-quotient (resolvent) axiom is the
single substantive demand; any finite-defect realization satisfying it is a de
Branges--Pontryagin space whose structure function shares its wrong-half-plane zeros with
$E_\xi$, the full identification $E=E_\xi$ remaining open. We close with the dichotomy that
explains why the spectral programmes converge to one wall: the symmetry data (spectrum,
anti-holomorphic real structure, resolvent) is unconditional, while the sign $\kappa=0$ is a
separate, global, one-sided ingredient living in the pair-correlation regime. Everything is
pure functional analysis; nothing here depends on, or proves, the Riemann Hypothesis.
\end{abstract}
\maketitle

\tableofcontents
\newpage

\section{Introduction}

Let $\xi(s)=\half s(s-1)\pi^{-s/2}\Gamma(s/2)\zeta(s)$ be the completed Riemann zeta
function, entire of order one, with $\xi(s)=\xi(1-s)$ and real on the critical line, and
with nontrivial zeros $\rho=\beta+i\gamma$ confined to the open strip $0<\beta<1$. The
Riemann Hypothesis ($\RH$) asserts $\beta=\half$ for all of them.

By Weil's explicit formula, $\RH$ is equivalent to the positivity of a single Hermitian
form. For a suitable even real test function $f$, the sum
\begin{equation}\label{eq:weil-id}
\sum_\rho\widehat f(\gamma_\rho)\;=\;A_\infty(f)-P(f),
\end{equation}
with $A_\infty$ the archimedean (gamma-factor) term and $P$ the sum over prime powers,
extends to a quadratic form $Q(f,f)=\sum_\rho\varphi_f(\rho)^2$ whose positivity for all $f$
is precisely $\RH$. The form $Q$ is \emph{indefinite}: it is a Kre\u\i n form, and its
negative index
\[
\kappa\;=\;\#\{\rho:\beta_\rho>\tfrac12\}\;=\;2m
\]
equals the number of off-line zeros lying in the right half-strip --- equivalently
$2m$, where $m$ is the number of Klein quartets $\{\half\pm b\pm i\gamma\}$ of off-line
zeros --- so that $\kappa=0\iff\RH$. (The negative index counts the zeros of the structure
function $E_\xi$ in one open half-plane, two per quartet, not all four off-line zeros; the
parity $\kappa=2m$ and its arithmetic consequences are the subject of the companion
paper~\cite{Defect}.) This reformulation is the
starting point of a long sequence of spectral attacks (Hilbert--P\'olya and its de Branges
canonical-system form, Connes' scaling flow, Burnol's de Branges spaces, and the conjectural
arithmetic surface for $\operatorname{Spec}\mathbb Z$): each seeks to \emph{realize} $Q$ as
the inner product of a Hilbert or Pontryagin space whose spectrum is $\{\gamma_\rho\}$ and
whose positivity would force the zeros onto the line.

The purpose of this paper is to give a complete, self-contained anatomy of $Q$ as an
indefinite form, and of the realization problem, organized around one localization. We prove,
from an autonomous definition of $Q$, that the entire indefiniteness is carried by the
\emph{off-axis evaluation functionals} attached to off-line zeros
(\S\ref{sec:setup}--\S\ref{sec:localization}); we then deploy this localization in three
directions:
\begin{itemize}
\item \textbf{Obstruction} (\S\ref{sec:modular}--\S\ref{sec:obstruction}). An entire class of
candidate realizations is excluded: positive-kernel reproducing spaces, fixed-sign Cauchy
pairings, and forms depending smoothly and non-degenerately on $\beta$. The modular surface
furnishes the sharpest concrete instance — its resonance pairings are sign-definite for every
zero configuration — and the three general theorems explain why. A genuine carrier must have
an \emph{indefinite kernel} or a \emph{degeneration at the line}.
\item \textbf{Mechanism and count} (\S\ref{sec:mechanism}--\S\ref{sec:regimes}). The Weil form
is bilinear, and that is exactly why it can be indefinite; the local mechanism is one
Lorentzian $(1,1)$ block per off-axis evaluation, and the global count
$\kappa=2m$ follows from an independence lemma together with the de
Branges Pontryagin-index theorem in the finite-index regime, whose scope we delimit.
\item \textbf{Rigidity} (\S\ref{sec:rigidity}). Among the three de Branges axioms, only the
difference-quotient (resolvent) axiom carries new content; any finite-defect realization
satisfying it is a de Branges--Pontryagin space whose structure function shares the
wrong-half-plane zeros of $E_\xi$. The residual problem is two explicit items: the resolvent
axiom and the identification $E=E_\xi$.
\end{itemize}
We close (\S\ref{sec:dichotomy}) with the dichotomy \emph{symmetry versus positivity}: the
symmetry data of any realization is unconditional, while $\kappa=0$ is orthogonal to it and
irreducibly global, which is why the spectral programmes all terminate at the same wall —
the Weil positivity in the pair-correlation regime.

Everything below is pure functional analysis. None of it assumes $\RH$, and none of it proves
$\RH$; the contribution is structural — a precise account of where the indefiniteness of $Q$
lives, which realizations are excluded, and what an admissible realization must look like. The
exact arithmetic count of negative directions per off-line orbit (the parity refinement
$\kappa=2m$) and the resulting modified explicit formula are taken up in a companion
paper~\cite{Defect}; here we keep to the de Branges count $\kappa=2m$ (zeros of $E_\xi$ in one open half-plane)
and its local mechanism.

\section{The Weil form, defined autonomously}\label{sec:setup}

We define $Q$ directly on a concrete test space, so that all later statements are about an
explicit object rather than about the bookkeeping of the explicit formula. This is the
viewpoint that makes the localization theorem a self-contained statement.

\begin{definition}[Test space and transform]\label{def:setup}
Let $\mathcal V$ be the space of even real $f\in C_c^\infty(\mathbb R)$. For $f\in\mathcal V$
define the entire transform
\[
\varphi_f(s)\;=\;\int_{\mathbb R} f(u)\,e^{(s-\frac12)u}\,du,\qquad s\in\mathbb C .
\]
Since $f$ has compact support, $\varphi_f$ is entire of exponential type. Writing
$s=\half+iz$ we have $\varphi_f(\half+iz)=\int f(u)e^{izu}\,du=\fhat(z)$, the Fourier
transform of $f$ extended to $z\in\mathbb C$; because $f$ is even and real, $\fhat$ is even
and real on $\mathbb R$.
\end{definition}

\begin{definition}[The localized Weil form]\label{def:weil}
For $f,g\in\mathcal V$ set
\[
Q(f,g)\;=\;\sum_{\rho}\varphi_f(\rho)\,\varphi_g(\rho),
\]
the sum running over the nontrivial zeros $\rho$ of $\zeta$ with multiplicity.
\end{definition}

\begin{lemma}[Absolute convergence and symmetries]\label{lem:conv}
For $f,g\in\mathcal V$ the series in Definition~\ref{def:weil} converges absolutely, and
\[
\varphi_f(1-s)=\varphi_f(s),\qquad \varphi_f(\bar s)=\overline{\varphi_f(s)} .
\]
\end{lemma}
\begin{proof}
Evenness of $f$ gives $\varphi_f(1-s)=\int f(u)e^{(\frac12-s)u}du=\int f(-u)e^{(\frac12-s)u}
du=\int f(u)e^{(s-\frac12)u}du=\varphi_f(s)$; reality of $f$ gives
$\overline{\varphi_f(s)}=\int f(u)e^{(\bar s-\frac12)u}du=\varphi_f(\bar s)$. For
convergence, write $\rho=\beta+i\gamma$. With $f$ supported in $[-T_0,T_0]$,
$|\varphi_f(\rho)|=|\fhat(\gamma-i(\beta-\half))|\le\|f\|_1\,e^{T_0|\beta-\frac12|}\le
\|f\|_1\,e^{T_0/2}$, while integrating by parts twice against $e^{(\rho-\frac12)u}$ yields
$|\varphi_f(\rho)|\le C_f\,(1+|\gamma|)^{-2}$ uniformly for $0\le\beta\le1$. The Riemann--von
Mangoldt density $\#\{\rho:|\gamma|\le T\}\sim\frac{T}{\pi}\log\frac{T}{2\pi}$ makes
$\sum_\rho(1+|\gamma|)^{-2}<\infty$, hence $\sum_\rho|\varphi_f(\rho)\varphi_g(\rho)|<\infty$.
\end{proof}

The two symmetries say that the summand $\varphi_f(\rho)^2$ is invariant under
$\rho\mapsto1-\rho$ and is conjugated under $\rho\mapsto\bar\rho$; this is what couples the
four members of an off-line orbit, computed next.

\section{Localization of the index in off-axis evaluation}\label{sec:localization}

The nontrivial zeros split into the on-line zeros $\rho=\half+i\gamma$ and the off-line
zeros, which by the symmetries of Lemma~\ref{lem:conv} occur in \emph{Klein quartets}
\[
\{\rho,\ \bar\rho,\ 1-\rho,\ 1-\bar\rho\}=\{\tfrac12\pm b\pm i\gamma\},\qquad
b=\beta-\tfrac12\neq0 .
\]

\begin{theorem}[Localization of the negative directions]\label{thm:loc}
On $\mathcal V$ the Weil form decomposes as $Q=Q_{\mathrm{on}}+Q_{\mathrm{off}}$, where:
\begin{enumerate}
\item each on-line zero $\half+i\gamma$ contributes the manifestly nonnegative term
$\varphi_f(\half+i\gamma)^2=\fhat(\gamma)^2\ge0$, with $\fhat(\gamma)\in\mathbb R$;
\item each off-line quartet $\{\half\pm b\pm i\gamma\}$ contributes, when restricted to even
real test functions, the indefinite Lorentzian block
\[
4\,\Re\!\big[\fhat(\gamma-ib)^2\big]\;=\;4\,(u^2-v^2),\qquad
u=\Re\fhat(\gamma-ib),\ \ v=\Im\fhat(\gamma-ib),
\]
built solely from the off-axis evaluation functional $E_{\gamma-ib}:f\mapsto\fhat(\gamma-ib)$.
\end{enumerate}
Hence the indefiniteness is carried entirely by the off-axis evaluations: $Q_{\mathrm{off}}
\equiv0\iff$ there are no off-line zeros $\iff\RH$. The on-axis functionals ($b=0$) are the
positive limit of the off-axis ones as $\beta\to\half$. The exact negative index is computed
in Corollary~\textnormal{\ref{cor:count}}.
\end{theorem}
\begin{proof}
For an on-line zero, $\varphi_f(\half+i\gamma)=\fhat(\gamma)$ is real (Lemma~\ref{lem:conv}),
so its square is $\ge0$. For an off-line quartet restricted to even real $f$, group the four
summands by the symmetries. Put $w=\fhat(\gamma-ib)=\varphi_f(\half+b+i\gamma)$. By
$\varphi_f(1-s)=\varphi_f(s)$ the zero $1-\rho=\half-b-i\gamma$ gives $w$, and by
$\varphi_f(\bar s)=\overline{\varphi_f(s)}$ the zeros $\bar\rho$ and $1-\bar\rho$ give
$\bar w$; the quartet contributes $w^2+w^2+\bar w^2+\bar w^2=4\Re[w^2]=4(u^2-v^2)$, the form
$\mathrm{diag}(4,-4)$ in $(u,v)$, depending on $f$ only through $w=E_{\gamma-ib}(f)$. This
exhibits the off-line indefiniteness on the real test space; the full (complex) negative
index is the de Branges Pontryagin index of the next subsection.
\end{proof}

\begin{remark}[Real-test-space mechanism versus the complex index]\label{rem:two-blocks}
On the \emph{even real} space $\mathcal V$ the reality $\fhat(\bar z)=\overline{\fhat(z)}$
forces the evaluations at the two heights $\pm\gamma$ of a quartet to be complex conjugates,
so a quartet manifests there as the \emph{single} Lorentzian block above. In the complex de
Branges space $\HE$, however, the evaluations at the two upper-half-plane zeros of $E_\xi$ in
a quartet are genuinely independent, and the quartet carries \emph{two} negative directions.
The resulting count is $\kappa=2m$ (Corollary~\ref{cor:count}); the passage from the
real-restricted mechanism to the complex parity is exactly the phenomenon analyzed
in~\cite{Defect}.
\end{remark}

\begin{remark}[Four load-bearing choices]\label{rem:choices}
The decomposition rests on exactly four explicit choices, each isolated above: the test space
$\mathcal V$ (even real, compactly supported), the transform $\varphi_f$ (which makes the
evaluation entire and ties $\half+iz$ to $\fhat(z)$), the quartet regrouping (which uses both
symmetries of Lemma~\ref{lem:conv}), and the convergence input (Riemann--von Mangoldt). None
is a corollary of explicit-formula bookkeeping; together they make Theorem~\ref{thm:loc} an
autonomous statement about the explicit object of Definitions~\ref{def:setup}--\ref{def:weil}.
\end{remark}

Theorem~\ref{thm:loc} reduces the realization problem to a single analytic object: a function
space on which the off-axis point evaluations $E_{\gamma-ib}$ ($b\ne0$) are bounded and carry
the indefinite combination $u^2-v^2$. On $L^2(\mathbb R)$ these evaluations are unbounded, so
$L^2$ is excluded at the outset; the question is which spaces support them, and whether
boundedness forces positivity. The next three sections answer this.

\section{A definite test case: the modular-surface resonances}\label{sec:modular}

Before the general obstruction we record the sharpest concrete failure, on the most natural
geometric home of the zeros. On the modular surface $X=SL_2(\mathbb Z)\backslash\mathbb H$ the
zeros appear as scattering resonances: the constant term of the Eisenstein series $E(z,s)$ is
$y^s+\varphi(s)y^{1-s}$ with scattering determinant
\[
\varphi(s)=\frac{\xi(2s-1)}{\xi(2s)},
\]
whose poles are exactly the resonances $s_\rho=\rho/2=\tfrac14+i\gamma_\rho/2$ (since
$\xi(2s)=0\iff2s=\rho$). One asks whether the resonance space carries the indefinite Weil
form. It does not, for a structural reason that the next section generalizes.

\begin{proposition}[The resonance weights are $O(1)$]\label{prop:weights}
The resonance state $\Psi_\rho=\Res_{s=s_\rho}E(\cdot,s)$ has cusp weight
$R_\rho=\Res_{s=\rho/2}\varphi(s)=\xi(\rho-1)/2\xi'(\rho)$, and $|R_\rho|=O(1)$: the
archimedean factors cancel, $\xi(\rho-1)\asymp e^{-\pi\gamma/4}\asymp\xi'(\rho)$, so the ratio
is bounded (numerically $0.52$--$0.66$ for the first fifteen zeros, slowly increasing). The
resonance space is therefore non-trivial, in contrast to the de Branges space $\HE$ of the
critical-line $\xi$, whose weights decay like $e^{-\pi t/4}$.
\end{proposition}

Thus the modular surface passes the spectrum test with genuine, bounded-below weights. The
pairing is where it fails.

\begin{theorem}[The resonance pairings are half-plane definite]\label{thm:modular}
Write $\rho=\beta_\rho+i\gamma_\rho$ with $0<\beta_\rho<1$ and $s_\rho=\rho/2$.
\begin{enumerate}
\item \emph{(Sesquilinear.)} The regularized Maass--Selberg--Zagier inner product of resonance
states has Gram matrix $G=DKD^{*}$, $D=\mathrm{diag}(R_\rho)$,
$K_{\rho\sigma}=(s_\rho+\overline{s_\sigma}-1)^{-1}$, and $K$ is negative-definite for every
configuration of zeros in $0<\beta<1$; hence $G$ is negative-definite, of constant signature.
\item \emph{(Bilinear.)} The bilinear Gamow/Jost self-pairing of $\rho$ with its conjugate is
\[
B(\Psi_\rho,\Psi_{\bar\rho})=\frac{-R_\rho R_{\bar\rho}}{s_\rho+s_{\bar\rho}-1}
=\frac{|R_\rho|^2}{1-\beta_\rho}>0
\]
for every $\beta_\rho\in(0,1)$, on-line and off-line alike.
\end{enumerate}
Consequently neither pairing can equal $Q$: both are sign-definite throughout the strip, so
the negative index $\kappa$ is structurally absent from them, with or without $\RH$.
\end{theorem}
\begin{proof}
(1) Set $w_\rho=\half-s_\rho$, so $\Re w_\rho=\tfrac{1-\beta_\rho}{2}>0$ because
$\beta_\rho<1$. Then
\[
K_{\rho\sigma}=\frac{1}{s_\rho+\overline{s_\sigma}-1}=\frac{-1}{w_\rho+\overline{w_\sigma}}
=-\int_0^\infty e^{-w_\rho t}\,\overline{e^{-w_\sigma t}}\,dt,
\]
the negative of the Cauchy/Szeg\H o reproducing kernel of the Hardy space $H^2$ of the right
half-plane, hence negative-definite; conjugating by the invertible $D$ preserves signature.
(2) Use $R_{\bar\rho}=\overline{R_\rho}$ and $s_\rho+s_{\bar\rho}-1=\beta_\rho-1$.
\end{proof}

\begin{remark}
The obstruction is that both Gram kernels are Cauchy-type $1/(s_\rho\pm\overline{s_\sigma}-c)$
with denominator of one sign throughout $0<\beta<1$, while $\kappa$ must detect $\beta=\half$
\emph{exactly} — a measure-zero, non-monotone condition. A pairing smooth and sign-stable
across the strip cannot witness it. Every natural pairing the surface supplies — the analytic
Maass--Selberg--Zagier product and the arithmetic Arakelov/N\'eron--Tate pairing
(negative-definite by Faltings--Hriljac) — is definite. This is the general phenomenon, to
which we now turn.
\end{remark}

\section{The obstruction theorems and the excluded class}\label{sec:obstruction}

We isolate three properties shared by the failing constructions and prove that each forces
definiteness or non-detection. The results are elementary; the value is the classification.

\subsection{Positive kernels are semidefinite}
Let $\Omega$ be a set and $K:\Omega\times\Omega\to\mathbb C$ a positive-definite kernel; let
$\HK$ be its Moore--Aronszajn reproducing-kernel Hilbert space, with $K_z=K(\cdot,z)$ and
$\langle f,K_z\rangle=f(z)$.

\begin{theorem}[Positive kernel $\Rightarrow$ semidefinite]\label{thm:A}
Any Hermitian form realized as $B(f,g)=\langle f,g\rangle_{\HK}$ for a positive-definite kernel
$K$ is positive semidefinite; its negative index is $0$. In particular, if $Q$ is realized as
the inner product of a positive-kernel RKHS, then $\kappa=0$, i.e.\ $\RH$ — the form never
exhibits a negative direction.
\end{theorem}
\begin{proof}
For $f=\sum_i c_iK_{z_i}$, $\|f\|_{\HK}^2=\sum_{i,j}c_i\bar c_j K(z_i,z_j)\ge0$ by
positive-definiteness; such $f$ are dense, so $B\succeq0$ and its negative index vanishes.
\end{proof}

This is the Moore--Aronszajn tautology read backwards: indefiniteness cannot be manufactured
from a positive kernel. Bergman, Hardy, Fock, Mehler/heat, Bochner-positive convolution
pairings, and standard de Branges spaces with a Hermite--Biehler structure function are all
excluded as carriers of $\kappa$.

\subsection{Fixed-sign Cauchy kernels are definite}
Theorem~\ref{thm:modular} is the special case $c=1$ of the following.

\begin{theorem}[Fixed-sign Cauchy kernel $\Rightarrow$ definite]\label{thm:B}
Let $c\in\mathbb R$ and $S\subset\mathbb C$ with all $\Re z<c/2$ (resp.\ all $\Re z>c/2$).
Then $K(z,w)=(z+\bar w-c)^{-1}$ is negative-definite (resp.\ positive-definite) on $S$, and
any Gram form $G=[d_i\bar d_j K(z_i,z_j)]$ with weights $d_i\ne0$ has constant signature.
\end{theorem}
\begin{proof}
With $w_z=\tfrac c2-z$ (so $\Re w_z>0$ on $S$ in the first case),
$K(z,w)=-(w_z+\overline{w_w})^{-1}=-\int_0^\infty e^{-w_zt}\overline{e^{-w_wt}}\,dt$ is the
negative of the half-plane Hardy kernel, hence negative-definite; the other case is identical
with $w_z=z-\tfrac c2$, and conjugation by $D=\mathrm{diag}(d_i)$ preserves signature.
\end{proof}

The hypothesis is exactly ``the parameters lie in a half-plane not straddling $\Re=c/2$.''
For the modular resonances $s_\rho=\rho/2$ this holds because $0<\beta_\rho<1$ confines them
to $\Re s_\rho<\half$; the bounded critical strip forces definiteness regardless of $\RH$.

\subsection{No $\beta$-smooth non-degenerate form detects the line}
The condition $\beta=\half$ has real codimension one; signature is locally constant on the
non-degenerate locus.

\begin{theorem}[Smoothness $\Rightarrow$ no detection]\label{thm:C}
Let $\beta\mapsto Q_\beta$ be a Hermitian form on a fixed finite-dimensional space (or with
discrete, continuously varying spectrum), depending continuously on $\beta$ near $\half$. If
$Q_\beta$ is non-degenerate throughout a neighborhood of $\half$, its negative index
$\kappa(\beta)$ is constant there. In particular $Q_\beta$ cannot have $\kappa(\half)=0$ and
$\kappa(\beta)\ge1$ for $\beta\ne\half$ unless it degenerates at $\half$.
\end{theorem}
\begin{proof}
Eigenvalues vary continuously; $\kappa(\beta)$, the number of negative ones, can change only
where an eigenvalue crosses $0$, i.e.\ at a degeneracy. On a non-degenerate neighborhood none
crosses, so $\kappa$ is constant. The stated detection pattern forces a zero eigenvalue at
$\half$.
\end{proof}

This explains the recurrent finding that every natural pairing tested was a smooth,
non-degenerate function of $\beta$ — the Petersson form $\sim(1-\beta)$, the bilinear Gamow
form $\sim1/(1-\beta)$, the Klein--Gordon current ($\beta$-independent) — none vanishing at
$\half$, hence all of constant signature and blind to $\RH$. The genuine detector, the off-line
displacement $b=\beta-\half$, \emph{does} vanish at $\half$: it is degenerate there by design.

\begin{corollary}[The excluded class and the two escape routes]\label{cor:class}
The localized Weil form $Q$ — indefinite, with $\kappa$ a genuine, $\RH$-detecting index —
cannot be realized by any of: \emph{(i)} a positive-definite reproducing kernel
(Theorem~\ref{thm:A}); \emph{(ii)} a Cauchy kernel $1/(z+\bar w-c)$ whose parameter domain
lies on one side of $\Re=c/2$, in particular any modular/scattering/Arakelov pairing of zeros
confined to the strip (Theorem~\ref{thm:B}); \emph{(iii)} a non-degenerate Hermitian form
depending smoothly on $\beta$ near $\half$ (Theorem~\ref{thm:C}). Consequently a genuine
realization must possess at least one feature its competitors lack: an \emph{indefinite
(Pontryagin) kernel}, or a \emph{degeneration at $\beta=\half$}.
\end{corollary}

The second escape route is exactly the off-axis evaluation of Theorem~\ref{thm:loc}, which
degenerates as $b\to0$. The first is the indefinite de Branges space, whose local mechanism we
now identify.

\section{The local mechanism and the exact count}\label{sec:mechanism}

\subsection{Bilinear, not sesquilinear}
For even real $f$, $\fhat$ is even and real on $\mathbb R$, so the entire extension of
$|\fhat|^2$ is the holomorphic square $\fhat(z)^2$, not $\fhat(z)\overline{\fhat(z)}$. Hence
\begin{equation}\label{eq:bilinear}
Q(f,g)=\sum_\rho\fhat(z_\rho)\,\ghat(z_\rho),\qquad z_\rho=\tfrac{\rho-\frac12}{i},
\end{equation}
is a \emph{symmetric bilinear} form, not the sesquilinear $\sum\fhat(z_\rho)\overline{\ghat
(z_\rho)}$. This distinction is the whole point: the bilinear square of a complex number is
indefinite.

\begin{proposition}[Sesquilinear vs.\ bilinear]\label{prop:vs}
Let $E_z(f)=\fhat(z)$. The sesquilinear form $f\mapsto|E_z(f)|^2$ is $\ge0$ for every $z$. The
bilinear form $f\mapsto E_z(f)^2$ is $\ge0$ where $E_z(f)$ is real (e.g.\ $z\in\mathbb R$) but
indefinite where $E_z(f)$ ranges over $\mathbb C$: with $E_z(f)=u+iv$, $\Re[E_z(f)^2]=u^2-v^2$.
\end{proposition}

\subsection{The dichotomy theorems}
\begin{theorem}[Sesquilinear evaluation forms are positive]\label{thm:E}
In any reproducing-kernel Hilbert space with continuous evaluations $E_z(f)=\langle f,K_z
\rangle$, and any points $z_i$ and weights $c_i\ge0$, the form $f\mapsto\sum_i c_i|E_{z_i}(f)|^2$
is positive semidefinite.
\end{theorem}
\begin{proof}
$|E_{z_i}(f)|^2=\langle f,(K_{z_i}\otimes K_{z_i})f\rangle$ with $K_{z_i}\otimes K_{z_i}
\succeq0$; a nonnegative combination of positive forms is positive.
\end{proof}

\begin{theorem}[Bilinear evaluation squares: signature $(1,1)$ off-axis]\label{thm:F}
The real quadratic form $f\mapsto\Re[E_z(f)^2]$ has signature $(1,0)$ where $E_z(f)$ is forced
real (on-axis) and signature $(1,1)$ where $E_z(f)=u+iv$ realizes all of $\mathbb C$
(off-axis). Thus each off-axis evaluation contributes \emph{at most} one negative direction.
\end{theorem}
\begin{proof}
Factor $E_z=u+iv$ with $u,v$ real-linear. On-axis $v\equiv0$ and $\Re[E_z^2]=u^2$: signature
$(1,0)$. Off-axis $(u,v)$ maps onto $\mathbb R^2$ and $\Re[E_z^2]=u^2-v^2$ pulls back
$\mathrm{diag}(1,-1)$: signature $(1,1)$.
\end{proof}

\subsection{From blocks to count}
The negative index of a \emph{sum} $\sum_\rho(u_\rho^2-v_\rho^2)$ is not automatically the
number of blocks: the $v_\rho$ could be dependent or interact with the $u_\rho$. The exact
count needs an independence input.

\begin{lemma}[Independence of off-axis evaluations]\label{lem:indep}
On $\mathcal V$ (indeed on $\mathcal S(\mathbb R)$), for distinct off-axis points
$z_1,\dots,z_n$ the $2n$ real-linear functionals $\{\Re E_{z_i},\Im E_{z_i}\}$ are linearly
independent; hence $\sum_{i=1}^n((\Re E_{z_i})^2-(\Im E_{z_i})^2)$ has negative index exactly
$n$.
\end{lemma}
\begin{proof}
A relation $\sum_i c_iE_{z_i}=0$ on $\mathcal S$ reads $\int f(x)\sum_i c_ie^{-iz_ix}\,dx=0$
for all $f\in\mathcal S$, forcing the tempered distribution $\sum_i c_ie^{-iz_ix}=0$;
exponentials of distinct frequencies are independent (nonvanishing Vandermonde), so all
$c_i=0$. Thus the $E_{z_i}$ are $\mathbb C$-independent, the $\{\Re E_{z_i},\Im E_{z_i}\}$ are
$\mathbb R$-independent, and in a dual basis the form is $\mathrm{diag}(1,-1,\dots,1,-1)$.
\end{proof}

\begin{corollary}[The count]\label{cor:count}
On the de Branges space $\HE$, $E_\xi(z)=\xi(\half+iz)$, the negative (Pontryagin) index of
$Q$ equals the number of zeros of $E_\xi$ in the open upper half-plane. A zero $\rho$ of
$\xi$ maps to $z=-i(\rho-\half)$ with $\Im z=\half-\Re\rho$, so the upper-half-plane zeros are
exactly those with $\Re\rho<\half$ --- two per off-line quartet. Hence
\[
\kappa\;=\;\#\{\rho:\Re\rho<\tfrac12\}\;=\;\#\{\rho:\Re\rho>\tfrac12\}\;=\;2m,
\]
$m$ the number of off-line quartets. For finitely many off-axis points the independence is
Lemma~\ref{lem:indep}; for the full zero set the count is the de Branges Pontryagin-index
theorem~\cite{deBranges}, of which Theorem~\ref{thm:F} and Lemma~\ref{lem:indep} are the local
mechanism. The parity $\kappa\equiv0\pmod2$ and the per-quartet structure are developed
in~\cite{Defect}.
\end{corollary}

\begin{remark}[Two reductions deferred]\label{rem:deferred}
Lemma~\ref{lem:indep} is the finite, complexified mechanism. Two further reductions — the
passage to the real even test space, where the conjugate pairing $z_i\leftrightarrow-\bar z_i$
couples the off-axis blocks, and the convergence of the full infinite zero sum — are supplied
by the de Branges Pontryagin theory. The first is also where the arithmetic refinement
enters: counting the negative directions \emph{per Klein orbit} (rather than per evaluation
point) gives the parity statement $\kappa=2m$ with $m$ the number of off-line orbits, proved
in the companion paper~\cite{Defect}. In the present functional-analytic count we use the de
Branges value $\kappa=2m$.
\end{remark}

\subsection{Signature from the spectrum}
\begin{proposition}[Definiteness is governed by real vs.\ complex spectral points]\label{prop:spectrum}
For a Weil-type bilinear form $\sum_\rho E_{z_\rho}^2$ on a de Branges space $H(E)$ (where the
$z_\rho$ are the zeros of $E$): each real $z_\rho$ contributes signature $(1,0)$ and each
non-real $z_\rho$ contributes $(1,1)$. Hence the form is positive-definite iff all $z_\rho$
are real, and has negative index $\#\{\text{non-real }z_\rho\}$ (modulo Lemma~\ref{lem:indep})
otherwise. Equivalently: $E$ Hermite--Biehler $\Leftrightarrow$ all zeros real
$\Leftrightarrow$ definite; $E$ non-Hermite--Biehler with $\kappa$ wrong-half zeros
$\Leftrightarrow$ indefinite of index $\kappa$.
\end{proposition}

\begin{center}
\renewcommand{\arraystretch}{1.2}\small
\begin{tabular}{|l|c|c|c|}
\hline
\textbf{Space} & \textbf{off-axis $E_z$} & \textbf{spectrum} & \textbf{bilinear signature}\\\hline
Hardy / Bergman / Fock / Paley--Wiener & bounded & real (HB) & definite\\
de Branges $H(E)$, $E$ Hermite--Biehler & bounded & real & definite\\
\textbf{Pontryagin--de Branges, $E$ non-HB} & bounded & $\kappa$ complex & \textbf{indefinite, index $\kappa$}\\
$L^2(\mathbb R)$ & \emph{unbounded} & --- & ---\\\hline
\end{tabular}
\end{center}

The indefiniteness is thus never an artifact of unbounded evaluation; it is the presence of
$\kappa$ non-real spectral points (off-line zeros), each a $(1,1)$ block by
Theorem~\ref{thm:F}. The indefinite de Branges space $\HE$ is the canonical known carrier; we
do not claim it is the unique one (a full classification of indefinite realizations is open).

\section{The three index regimes and the finiteness question}\label{sec:regimes}

The count of Corollary~\ref{cor:count} invokes the de Branges Pontryagin-index theorem, which
presupposes a \emph{finite} negative index. We make this scope precise.

\begin{definition}
A \emph{Pontryagin space} $\Pi_\kappa$ has finite negative index $\kappa<\infty$; a
\emph{Kre\u\i n space} allows $\kappa=\infty$. A de Branges space $H(E)$ is a Hilbert space
when $E$ is Hermite--Biehler, and a $\Pi_\kappa$ when $E$ has exactly $\kappa<\infty$ zeros in
the open upper half-plane; for infinitely many such zeros its negative index is infinite.
\end{definition}

\begin{proposition}[Three regimes]\label{prop:regimes}
According to $\kappa=2m=\#\{\rho:\Re\rho>\tfrac12\}$, the realization $\HE$ of $Q$ falls into:
\[
\begin{array}{lll}
\text{(H)} & \kappa=0\ (\RH): & \text{a Hilbert space};\ E_\xi\ \text{Hermite--Biehler}.\\
\text{(P)} & 0<\kappa<\infty: & \text{a Pontryagin space}\ \Pi_\kappa.\\
\text{(K)} & \kappa=\infty: & \text{infinite index; the finite-index theorem does not apply}.
\end{array}
\]
The de Branges Pontryagin-index theorem, and the count of Corollary~\ref{cor:count}, are
available in cases \textnormal{(H)} and \textnormal{(P)} but not in \textnormal{(K)}.
\end{proposition}

\begin{proposition}[Where finiteness sits]\label{prop:finiteness}
The statement $\kappa<\infty$ (finitely many off-line zeros) is strictly implied by $\RH$
($\kappa=0$), hence a priori weaker than $\RH$; yet it is not implied by the known
unconditional results. The positive-proportion theorems (Levinson; Conrey) give $\ge2/5$ of
zeros on the line but bound nothing off it, and the zero-density estimates $N(\sigma,T)=
O(T^{c(1-\sigma)+\varepsilon})$ give $o(T)$ off-line zeros up to height $T$ for each fixed
$\sigma>\half$ but permit infinitely many in total. The only general constraint is that
off-line zeros occur in quartets contributing two negative directions each, so $\kappa\in2\mathbb Z_{\ge0}$ with minimal nonzero value $2$; this does not bound
$\kappa$, and both $\kappa<\infty$ and $\kappa=\infty$ remain open. We make no conjecture.
\end{proposition}

\begin{remark}
Thus the finite-vs-infinite index distinction is not a soft-analysis gap; it is the
zero-density frontier in index-theoretic dress. The de Branges realization is, strictly, a
finite-index analysis, clean under $\RH$ (case (H), the Hilbert case) and valid in (P); case
(K) would require a broader Kre\u\i n-space theory, which we do not develop. The quartet
constraint $\kappa\in2\mathbb Z_{\ge0}$ here is the coarse functional-analytic bound; the
arithmetic refinement to $\kappa=2m$ per orbit is in~\cite{Defect}.
\end{remark}

\section{Structural rigidity of the finite-defect realization}\label{sec:rigidity}

Existence of a realization does not entail uniqueness. We ask how rigid the de Branges
realization is, using de Branges' axiomatic characterization, extended to Pontryagin spaces by
Kre\u\i n and Langer.

\begin{definition}[de Branges axioms]\label{def:axioms}
Let $\Pi$ be a Pontryagin space of entire functions with indefinite inner product
$[\cdot,\cdot]$; write $\mathcal E_w:f\mapsto f(w)$ and $f^*(z)=\overline{f(\bar z)}$. Then
$\Pi$ is a \emph{de Branges--Pontryagin space} if:
\begin{itemize}
\item[\textbf{(dB1)}] $\mathcal E_w$ is continuous for every non-real $w$;
\item[\textbf{(dB2)}] \emph{(difference-quotient axiom)} if $f\in\Pi$ and $f(w)=0$ then
$z\mapsto f(z)/(z-w)\in\Pi$, and for $f,g\in\Pi$ with $f(w)=g(v)=0$,
\[
\Big[\tfrac{f(z)}{z-w},g\Big]-\Big[f,\tfrac{g(z)}{z-v}\Big]
=(\bar v-w)\Big[\tfrac{f(z)}{z-w},\tfrac{g(z)}{z-v}\Big];
\]
\item[\textbf{(dB3)}] \emph{(conjugation isometry)} $f\mapsto f^*$ maps $\Pi$ onto $\Pi$ with
$[f^*,g^*]=[g,f]$.
\end{itemize}
\end{definition}

\begin{theorem}[de Branges; Kre\u\i n--Langer~\cite{deBranges,KL}]\label{thm:dB}
A Pontryagin space of entire functions satisfying \textnormal{(dB1)--(dB3)} and containing a
nonzero element equals isometrically a space $H(E)$ for a structure function $E$ that is
Hermite--Biehler up to finitely many wrong-half-plane zeros, their number equal to the
negative index; $E$ is determined up to the real-unimodular normalization.
\end{theorem}

\begin{proposition}[Cost of each axiom]\label{prop:auto}
Let $H$ be a reproducing-kernel Pontryagin realization of $Q$ with finite $\kappa$, built from
the off-axis functionals of Theorem~\ref{thm:loc}. Then:
\begin{enumerate}
\item \textnormal{(dB3)} is automatic: by Lemma~\ref{lem:conv}, $\varphi_f(\bar s)=
\overline{\varphi_f(s)}$ and $\varphi_f(1-s)=\varphi_f(s)$, so $Q(f,g)=\sum_\rho\varphi_f(\rho)
\varphi_g(\rho)$ is invariant under $f\mapsto f^*$; hence $f\mapsto f^*$ is a
$[\cdot,\cdot]$-isometry.
\item \textnormal{(dB1)} is \emph{not} automatic: Theorem~\ref{thm:loc} bounds only the
specific functionals $E_{\gamma-ib}$ at the zero ordinates, whereas \textnormal{(dB1)} demands
continuity at \emph{every} non-real $w$ — the statement that $H$ is a reproducing-kernel space.
We take it as the standing hypothesis ``$H$ is a reproducing-kernel realization.''
\item \textnormal{(dB2)} is the single substantive demand: it asserts that $H$ is closed under
the resolvent-type operation $f\mapsto f/(z-w)$ obeying the de Branges commutation relation — a
shift/resolvent structure, the Hilbert--P\'olya object isolated as one hypothesis. It cannot be
read off from the symmetries of $Q$.
\end{enumerate}
\end{proposition}
\begin{proof}
(1) The two symmetries make $f^*$ a representative with conjugate $\varphi$-values at the
conjugate (and reflected) zeros; conjugation-symmetry of the zero set preserves the form value.
(2)--(3) A finite set of bounded functionals does not entail continuity at all non-real points;
and the difference-quotient identity is an operator-theoretic statement independent of the
symmetries.
\end{proof}

\begin{theorem}[Structural rigidity]\label{thm:rigid}
Let $H$ be a reproducing-kernel Pontryagin realization of $Q$ with finite $\kappa$
\textnormal{(}so \textnormal{(dB1)} holds\textnormal{)}, built from the off-axis functionals
of Theorem~\ref{thm:loc}, and suppose $H$ satisfies \textnormal{(dB2)}. Then:
\begin{enumerate}
\item $H=H(E)$ isometrically for some structure function $E$ with exactly $\kappa$ zeros in
the open upper half-plane;
\item the wrong-half-plane zeros of $E$ are precisely the off-line zeros of $\xi$ (in the
$z$-variable); so $E$ and $E_\xi(z)=\xi(\half+iz)$ share their upper-half-plane zero sets.
\end{enumerate}
\end{theorem}
\begin{proof}
By Proposition~\ref{prop:auto}, (dB3) holds and (dB1) is the standing hypothesis; (dB2) is
assumed. Theorem~\ref{thm:dB} gives $H=H(E)$ isometrically with negative index $\kappa$ equal
to the number of upper-half-plane zeros of $E$. By Theorem~\ref{thm:loc} and
Corollary~\ref{cor:count} the negative directions are indexed by the off-line zeros of $\xi$;
the isometry preserves the index and its carrier, so the wrong-half-plane zeros of $E$ are
exactly those off-line zeros, which are also the wrong-half-plane zeros of $E_\xi$.
\end{proof}

\begin{remark}[The identification $E=E_\xi$ is not established]\label{rem:identify}
Theorem~\ref{thm:rigid} pins the carrier to the de Branges--Pontryagin class and fixes its
wrong-half-plane zeros, but does not prove $E=E_\xi$: the structure function is underdetermined
by its upper-half-plane zeros alone, and $E$ could differ from $E_\xi$ by a Hermite--Biehler
factor with only real zeros (the on-line spectrum) and by the real-unimodular normalization.
A full identification would require matching the entire reproducing kernel of $Q$ to that of
$\xi$, including the on-line spectrum and boundary phase. We isolate this as the precise
remaining step; it is genuinely deeper than the class membership proved above.
\end{remark}

\begin{corollary}[The two residues]\label{cor:freedom}
In the finite-defect regime the realization problem splits into exactly two residues:
\textnormal{(i)} the resolvent axiom \textnormal{(dB2)}, and \textnormal{(ii)} the
identification $E=E_\xi$. Neither is supplied by the symmetries of $Q$. In the infinite-index
regime \textnormal{(K)} the finite-index theorem is silent and these tools do not decide the
realization.
\end{corollary}

\section{Symmetry versus positivity: why the programmes meet one wall}\label{sec:dichotomy}

We close with the structural dichotomy that organizes the entire analysis. Write
$\sigma(s)=\bar s$, $\iota(s)=1-s$, and $j=\sigma\iota:\,s\mapsto1-\bar s$, an
anti-holomorphic involution with $\Fix(j)=\{\Re s=\half\}$.

\begin{definition}[Realization and its symmetry data]\label{def:real}
A \emph{realization} of $Q$ is a reproducing-kernel (Pontryagin or Kre\u\i n) space on which
$Q$ acts as the inner product through the off-axis functionals of Theorem~\ref{thm:loc}. Its
\emph{symmetry data} is the triple (spectrum $\{\gamma_\rho\}$; real structure $j$;
de Branges resolvent (dB2)). Its \emph{positivity} is the scalar $\kappa$, with
$\RH\iff\kappa=0$.
\end{definition}

\begin{proposition}[Symmetry is unconditional]\label{prop:sym}
$\sigma,\iota$ generate a Klein four-group $V\cong(\mathbb Z/2)^2$ acting on $\mathbb C$, with
$j$ anti-holomorphic and $\Fix(j)$ the critical line; the zero set of $\xi$ is $V$-invariant
unconditionally. Consequently, for every admissible zero configuration — on-line or off —
there is a realization carrying the spectrum, the real structure $j$, and the resolvent: the
symmetry data is present whether or not $\RH$ holds.
\end{proposition}
\begin{proof}
$\sigma,\iota$ are commuting involutions; $j^2=\mathrm{id}$, $1-\bar s=s\iff\Re s=\half$. The
functional equation and real coefficients give $V$-invariance of the zeros. For an off-line
configuration the de Branges space of the corresponding $E$ (Hermite--Biehler up to finite
defect) carries multiplication (spectrum $=$ zeros), the conjugation $j$, and the
difference-quotient resolvent exactly as in the on-line case (Theorem~\ref{thm:rigid}); none
of this constrains $\kappa$.
\end{proof}

\begin{proposition}[The index is orthogonal to the symmetry]\label{prop:kappa}
The real structure acts on the spectrum setwise and pointwise:
\[
\text{(setwise)}\ \ j(\{\gamma_\rho\})=\{\gamma_\rho\};\qquad
\text{(pointwise)}\ \ j(\rho)=\rho\ \forall\rho .
\]
The setwise statement holds unconditionally (off-line zeros form $j$-symmetric quartets); the
pointwise statement is $\{\rho\}\subset\Fix(j)$, i.e.\ $\RH$. The index
$\kappa=\#\{\rho:\Re\rho>\tfrac12\}$ is the negative index of the bilinear pairing
(Theorem~\ref{thm:loc}); it is invariant under $V$ and commutes with the resolvent, hence
invisible to the symmetry data.
\end{proposition}
\begin{proof}
Off-line $\rho$ gives the $j$-symmetric orbit $\{\half\pm b\pm i\gamma\}$, so $j$ permutes
$\{\gamma_\rho\}$ regardless of $\RH$; $j$ fixes $\rho$ iff $\Re\rho=\half$, so pointwise
invariance of the whole spectrum is $\RH$. By Theorem~\ref{thm:loc} and
Corollary~\ref{cor:count}, $\kappa=\#\{\rho:\Re\rho>\tfrac12\}$ is unchanged by relabelling under
$V$ and by the resolvent, which preserve the inner product.
\end{proof}

\begin{proposition}[The negative directions are locally invisible]\label{prop:local}
Each negative direction of $Q$ is a frequency packet centred at an off-line ordinate
$\gamma_0$ of width below the mean zero spacing; it is orthogonal to the on-line evaluation
modes and is not lifted by any finite collection of neighbouring on-line zeros. Hence
$\kappa=0$ is implied by no local-in-$\gamma$ condition: the only constraint is the global
identity $\sum_\rho\varphi_f(\rho)^2=A_\infty(f)-P(f)$, and $\kappa=0\iff A_\infty\ge P$ on
all $f$ — the Weil criterion, a global positivity in the pair-correlation regime.
\end{proposition}
\begin{proof}[Proof sketch]
On an off-line quartet the form is $4(u^2-v^2)$, $u+iv=\fhat(\gamma_0-ib)$
(Theorem~\ref{thm:loc}); the negative direction suppresses $u$, maximizes $v$, and its
transform concentrates at $\gamma_0$ with frequency width controlled by the support of $f$,
which may exceed the inverse mean spacing — giving a packet narrower than the spacing. Such a
packet is orthogonal to the neighbouring on-line modes $\cos(\gamma't)$ and slips between the
discrete zeros, while the smooth archimedean density $A_\infty(f)=\int|\fhat|^2\Omega$,
$\Omega(r)\sim\log r$, still detects it; the identity routes the negativity to $P(f)>A_\infty
(f)$. Protection is the global comparison $A_\infty\ge P$, at the pair-correlation scale.
\end{proof}

\begin{theorem}[Symmetry versus positivity]\label{thm:dichotomy}
Let $R$ be any realization of $Q$ carrying the symmetry data of Definition~\ref{def:real}.
Then \emph{(i)} the symmetry data is $\RH$-independent (Proposition~\ref{prop:sym});
\emph{(ii)} $\RH$ is the positivity $\kappa=0$, the pointwise condition $\{\rho\}\subset
\Fix(j)$, neither implied by nor reflected in the symmetry data (Proposition~\ref{prop:kappa});
\emph{(iii)} the witnesses of $\kappa>0$ are sub-spacing and locally invisible, so $\kappa=0$
is irreducibly global — the Weil criterion $A_\infty\ge P$ in the pair-correlation regime
(Proposition~\ref{prop:local}). Hence no realization decides $\RH$ from its spectrum, real
structure, or resolvent: the sign is a separate, global, one-sided ingredient.
\end{theorem}
\begin{proof}
Immediate from Propositions~\ref{prop:sym}, \ref{prop:kappa}, \ref{prop:local}.
\end{proof}

\begin{corollary}[Convergence to one wall]\label{cor:converge}
The Hilbert--P\'olya / de Branges canonical system, Connes' scaling flow, Burnol's de Branges
spaces, and the conjectural arithmetic surface each supply the symmetry data and therefore, by
Theorem~\ref{thm:dichotomy}, terminate at the same residual statement: the global positivity
$\kappa=0$ in the pair-correlation regime. They differ in how they package the symmetry, not
in what they leave open.
\end{corollary}

\section{Conclusion}

The indefiniteness of the localized Weil form has a single, explicit source — the off-axis
evaluation functionals attached to off-line zeros (Theorem~\ref{thm:loc}) — and this one fact
organizes the whole landscape. It excludes an entire class of realizations (positive kernels,
fixed-sign Cauchy pairings, $\beta$-smooth forms; Corollary~\ref{cor:class}), with the modular
surface as the sharpest concrete casualty (Theorem~\ref{thm:modular}). It identifies the local
mechanism by which a Pontryagin index can nonetheless appear — one Lorentzian $(1,1)$ block per
off-axis evaluation (Theorem~\ref{thm:F}) — and reduces the global count to the de Branges
Pontryagin-index theorem in the finite-index regime (Corollary~\ref{cor:count},
\S\ref{sec:regimes}). It makes the de Branges realization structurally rigid up to two named
residues, the resolvent axiom and the identification $E=E_\xi$ (Theorem~\ref{thm:rigid},
Corollary~\ref{cor:freedom}). And it explains, through the dichotomy symmetry-versus-positivity
(Theorem~\ref{thm:dichotomy}), why the spectral programmes converge to one wall: the symmetry
data is abundant and unconditional, while the sign $\kappa=0$ is scarce, global, and one-sided.

The diagnostic conclusion is concrete. A successful realization must supply not a better
spectrum, real structure, or resolvent — these are RH-independent — but a genuinely new
\emph{global positivity}, carried by an indefinite kernel with bounded off-axis evaluation, or
equivalently by an indefinite arithmetic intersection that degenerates exactly at the critical
line. Nothing here proves the Riemann Hypothesis; it delimits, in pure functional-analytic
terms, what proving it through an operator realization would require.

\begin{thebibliography}{99}
\bibitem{Defect} D.~A.~Trejo Pizzo, \emph{The arithmetic defect of off-line zeros: parity, the
modified explicit formula, and two-scale structure}, companion paper, preprint, 2026.
\bibitem{deBranges} L.~de Branges, \emph{Hilbert Spaces of Entire Functions}, Prentice--Hall,
Englewood Cliffs NJ, 1968.
\bibitem{KL} M.~G.~Kre\u\i n and H.~Langer, \emph{\"Uber die verallgemeinerten Resolventen und
die charakteristische Funktion eines isometrischen Operators im Raume $\Pi_\kappa$}, in:
Hilbert Space Operators and Operator Algebras, Colloq.\ Math.\ Soc.\ J\'anos Bolyai
\textbf{5} (1972), 353--399.
\bibitem{Aronszajn} N.~Aronszajn, \emph{Theory of reproducing kernels}, Trans.\ Amer.\ Math.\
Soc.\ \textbf{68} (1950), 337--404.
\bibitem{Weil} A.~Weil, \emph{Sur les ``formules explicites'' de la th\'eorie des nombres
premiers}, Comm.\ S\'em.\ Math.\ Univ.\ Lund (1952), 252--265.
\bibitem{Iwaniec} H.~Iwaniec, \emph{Spectral Methods of Automorphic Forms}, Grad.\ Stud.\
Math.\ \textbf{53}, Amer.\ Math.\ Soc., 2002.
\bibitem{Zagier} D.~Zagier, \emph{The Rankin--Selberg method for automorphic functions which
are not of rapid decay}, J.\ Fac.\ Sci.\ Univ.\ Tokyo \textbf{28} (1981), 415--437.
\bibitem{Montgomery} H.~L.~Montgomery, \emph{The pair correlation of zeros of the zeta
function}, Proc.\ Sympos.\ Pure Math.\ \textbf{24} (1973), 181--193.
\bibitem{bgConnes99} A.~Connes, \emph{Trace formula in noncommutative geometry and the zeros of the Riemann zeta function}, Selecta Math.\ (N.S.) \textbf{5} (1999), 29--106.
\bibitem{bgCC21} A.~Connes and C.~Consani, \emph{Weil positivity and trace formula, the archimedean place}, Selecta Math.\ (N.S.) \textbf{27} (2021), no.~77.
\bibitem{bgBognar} J.~Bogn\'ar, \emph{Indefinite Inner Product Spaces}, Springer, 1974.
\bibitem{bgAzizov} T.~Ya.~Azizov and I.~S.~Iokhvidov, \emph{Linear Operators in Spaces with an Indefinite Metric}, Wiley, 1989.
\bibitem{bgBK} M.~V.~Berry and J.~P.~Keating, \emph{The Riemann zeros and eigenvalue asymptotics}, SIAM Rev.\ \textbf{41} (1999), 236--266.
\end{thebibliography}

\end{document}
